\section{Present Conservation, Creation as Relation, and Innermost Presence}

When people hear the word ``creation,'' they usually picture an initial spark at the dawn of time, after which the universe continues on its own momentum while its Maker watches from afar. That picture---the deist clockmaker who winds the spring and steps back---is the exact opposite of both classical metaphysics and Catholic doctrine. Once one distinguishes a cause of becoming from a cause of being, conservation in existence turns out not to be an optional extra that God performs after creating, but the uninterrupted continuation of the creative act itself.

\subsection{The Builder, the Sun, and Continuous Conservation}

Aquinas faces this question directly in \emph{ST} I, q.\ 104, a.\ 1: whether creatures need to be kept in being by God. He first sets aside an indirect sense of preservation---guarding a child from falling into fire, or shielding a metal from rust by removing corrupting agents. God preserves some creatures that way, yet incorruptible spiritual creatures like angels have no internal principle of material decay and do not need shielding from rust. The decisive question is \emph{per se} and direct preservation: whether a creature depends upon its preserver so radically that without that ongoing causal act it could not exist at all.

To answer, Aquinas sharpens the contrast between a cause of becoming and a cause of being. An agent that produces an effect of its own species---fire heating water, a human father begetting a child, a builder assembling stone and timber into a house---causes only the \emph{becoming} of the effect (\emph{causa fiendi}), not its \emph{being} (\emph{causa essendi}). The builder arranges pre-existing materials whose own natural forms keep the house standing after he walks home. Even hot water retains its warmth for a while after the fire is removed, because water is capable of receiving heat in the same specific way that heat exists in fire, however imperfectly.

By contrast, some effects cannot possess a perfection on the same footing as their source. Aquinas chooses his most famous analogy, drawn from pre-modern optics yet philosophically transparent in every age: sunlight illuminating the air. Air does not own light as an intrinsic property rooted in its own nature; it shines only while the sun actively illuminates it. ``The air does not continue to be lit up, even for a moment, when the sun ceases to act upon it.'' Turn off the sun, and darkness does not slowly creep in like cooling water; illumination vanishes instantaneously, because light has no independent root in the air.

Aquinas applies that analogy to every created thing, from the lowest pebble to the highest seraph:

\begin{quote}
Now every creature may be compared to God, as the air is to the sun which enlightens it. For as the sun possesses light by its nature, and as the air is enlightened by sharing the sun's nature; so God alone is Being in virtue of His own Essence, since His Essence is His existence; whereas every creature has being by participation, so that its essence is not its existence.\endnote{Aquinas, \emph{ST} I, q.\ 104, a.\ 1, corpus. Earlier in the same body he states the thesis directly: ``For the being of every creature depends on God, so that not for a moment could it subsist, but would fall into nothingness were it not kept in being by the operation of the Divine power.''}
\end{quote}

The Replies to Objections forestall two recurring misunderstandings. First, could an omnipotent God not simply \emph{grant} a creature the power to stay in existence on its own once created, just as a mortal builder leaves a house standing? Aquinas answers flatly no (\emph{ST} I, q.\ 104, a.\ 1, ad 2): ``God cannot grant to a creature to be preserved in being after the cessation of the Divine influence: as neither can He make it not to have received its being from Himself.'' That is no limit on divine omnipotence; it is simply the law of non-contradiction. A self-existing creature is a square circle. What exists by participation cannot simultaneously exist by essence.

Second, does conservation mean that God is constantly ``re-creating'' the world from scratch in flickering cinematographic frames, adding something new to the creature at every tick of the clock? Again Aquinas answers no (\emph{ST} I, q.\ 104, a.\ 1, ad 4): ``The preservation of things by God is a continuation of that action whereby He gives existence, which action is without either motion or time; so also the preservation of light in the air is by the continual influence of the sun.'' Conservation is not a series of discrete rescues; it is one timeless, continuous gift of being sustaining temporal creatures across every moment of their duration.

\subsection{Creation as Relation and God Innermost in All Things}

Two further metaphysical consequences follow directly from continuous conservation.

First, creation \emph{ex nihilo} is not a change or physical process (\emph{ST} I, q.\ 45, a.\ 2, ad 2). Every change presupposes a subject that was one way before and becomes another way afterward---marble uncarved becoming marble carved, cold water becoming hot. Creation produces the entire substance of a thing, presupposing nothing at all. When motion and successive transformation are stripped away from action and passion, what actually exists in the creature as ``being created''? Aquinas answers in \emph{ST} I, q.\ 45, a.\ 3: ``Creation places something in the thing created according to relation only\dots\ Hence creation in the creature is only a certain relation to the Creator as to the principle of its being.''\endnote{Aquinas, \emph{ST} I, q.\ 45, a.\ 3, corpus. In Reply~3, he clarifies that ``creation'' specifically imports this real relation of dependence \emph{with newness of being} (beginning after non-existence), whereas conservation designates the same ongoing relation of dependence across the creature's entire endurance.} To be a creature is not merely to have an interesting past; it is to be constituted, at every instant, by a real relation of total ontological dependence upon the Creator.

Second, because God is the immediate cause of \emph{esse} in every creature, God cannot act ``at a distance'' from across the universe. In \emph{ST} I, q.\ 8, a.\ 1, Aquinas explains that an agent must be immediately joined to that upon which it acts by its power. Because God is being itself by his essence, created being is his proper effect, just as igniting is the proper effect of fire---and he causes this effect not only when things begin, but ``as long as they are preserved in being.'' From this he draws his conclusion on divine presence:

\begin{quote}
Therefore as long as a thing has being, God must be present to it, according to its mode of being. But being is innermost in each thing and most fundamentally inherent in all things since it is formal in respect of everything found in a thing\dots\ Hence it must be that God is in all things, and innermostly.\endnote{Aquinas, \emph{ST} I, q.\ 8, a.\ 1, corpus. In Reply~1 and Reply~2, Aquinas balances immanence with transcendence: God is above all things by the excellence of his nature, and rather than being spatially contained by things, he spiritually contains the very things in which he is present.}
\end{quote}

Aquinas states the mode of this presence with care: God is in all things ``not, indeed, as part of their essence, nor as an accident, but as an agent is present to that upon which it works.'' That single sentence excludes pantheism without falling into deism. God is not blended into the creature's substance as a hidden ingredient, nor is he an external bystander peering in from beyond the stars. Because a creature's very act of existing is more intimate to it than its own qualities, thoughts, or choices---since none of those can even be unless it exists---the One who immediately bestows its \emph{esse} is present at a depth of its own reality that it cannot itself reach.

\subsection{Scripture and the Living God Who Upholds All Things}

Natural philosophy reaches this threshold by its own light: from change, compositeness, and the distinction of essence and existence to an uncreated, purely actual First Cause who continuously conserves every finite thing from within. Yet Christian thought did not invent this metaphysics in a vacuum, nor does Catholic faith rest on philosophical demonstration alone. What reason discerns under the concept of \emph{causa essendi} meets, is purified by, and stands under the revelation of the living God proclaimed in Sacred Scripture, expounded across the centuries by the Fathers of the Church, and solemnly confessed by the ecumenical councils.

Sacred Scripture voices this dependence in law, wisdom, psalm, and apostolic preaching. At the burning bush, God discloses his name to Moses as sheer, unborrowed being: ``I AM WHO AM\dots\ HE WHO IS, hath sent me to you'' (Exod.\ 3:14).\endnote{Scripture quotations in English follow the public-domain Challoner revision of the Douay-Rheims Bible; Latin checks follow the Clementine Vulgate (here \emph{Ego sum qui sum\dots\ Qui est, misit me ad vos}).} The Psalmist turns that truth straight into confession---``Know ye that the Lord he is God: he made us, and not we ourselves'' (Ps.\ 99:3)---and links every creature's breath directly to the Lord's unceasing regard: ``But if thou turnest away thy face, they shall be troubled: thou shalt take away their breath, and they shall fail, and shall return to their dust. Thou shalt send forth thy spirit, and they shall be created: and thou shalt renew the face of the earth'' (Ps.\ 103:29--30). The Book of Wisdom asks the exact question of continuous conservation: ``And how could any thing endure, if thou wouldst not? or be preserved, if not called by thee?'' (Wis.\ 11:26). In the New Testament, St.\ Paul tells the philosophers at the Areopagus that God is ``not far from every one of us. For in him we live and move and are'' (Acts 17:27--28); the Epistle to the Colossians declares of the Son that ``by him all things consist'' (Col.\ 1:17); and the Epistle to the Hebrews speaks of him ``upholding all things by the word of his power'' (Heb.\ 1:3), the exact text Aquinas places in the \emph{Sed contra} of \emph{ST} I, q.\ 104, a.\ 1.

Within that revealed horizon, Aquinas brings metaphysics and Exodus~3:14 directly together in \emph{ST} I, q.\ 13, a.\ 11, arguing that \emph{Qui Est}---``He Who Is''---is the most proper name of God among all names drawn from creaturely perfections. First, by its signification: every other name denotes some determinate form, whereas \emph{He Who Is} signifies existence itself, which in God alone is identical with his essence. Second, by its universality: every more specific name narrows and determines a mode of substance, falling short of God's infinity; hence Aquinas invokes St.\ John of Damascus's teaching in \emph{De Fide Orthodoxa} I.9 that this name denotes ``an infinite and indeterminate sea of substance''---a phrase Salmond renders directly from the Greek as ``a sea of essence infinite and unseen.'' Third, by its present tense: it signifies sheer present actuality, knowing neither past nor future.\endnote{Aquinas, \emph{ST} I, q.\ 13, a.\ 11, corpus (English Dominican rendering of Aquinas's Latin quotation of Damascene); compare St.\ John of Damascus, \emph{De Fide Orthodoxa} I.9 (S.\ D.\ F.\ Salmond translation, NPNF ser.\ 2, vol.\ 9, Edinburgh, 1898: ``For He keeps all being in His own embrace, like a sea of essence infinite and unseen''). In Reply~1, Aquinas adds that as regards the \emph{object} intended by the name (rather than its mode of signification), ``God'' is more proper as signifying the divine nature, and still more proper is the incommunicable Tetragrammaton signifying the singular substance of God himself.}

\subsection{Irenaeus, Athanasius, and Augustine on Sustained Being}

Long before scholastic terminology existed, the Fathers articulated the radical dependence of created being in response to Gnostic dualism and pagan theories of eternal matter.

In the late second century, St.\ Irenaeus of Lyons confronted critics who argued that if the human soul had a beginning in time, it must naturally perish like the body, whereas immortality could belong only to what was unbegotten. In \emph{Adversus Haereses} II.34.3--4, Irenaeus overturns that Greek premise at its root. Only God is unbegotten and immutable; all created things---stars, firmament, souls, and spirits alike---had a beginning and endure solely because God wills their continuance:

\begin{quote}
For life does not arise from us, nor from our own nature; but it is bestowed according to the grace of God\dots\ But as the animal body is certainly not itself the soul, yet has fellowship with the soul as long as God pleases; so the soul herself is not life, but partakes in that life bestowed upon her by God\dots\ When God therefore bestows life and perpetual duration, it comes to pass that even souls which did not previously exist should henceforth endure [for ever], since God has both willed that they should exist, and should continue in existence.\endnote{St.\ Irenaeus of Lyons, \emph{Adversus Haereses} II.34.3--4 (Ante-Nicene Fathers, vol.\ 1, Roberts and Rambaut translation).}
\end{quote}

A century and a half later, St.\ Athanasius of Alexandria placed continuous conservation at the center of his treatise \emph{Contra Gentes}. In chapters~41 and~42, Athanasius explains why the Word of God did not merely fashion the cosmos at the start and walk away:

\begin{quote}
For the nature of created things, inasmuch as it is brought into being out of nothing, is of a fleeting sort, and weak and mortal, if composed of itself only. But the God of all is good and exceeding noble by nature,---and therefore is kind. For one that is good can grudge nothing: for which reason he does not grudge even existence, but desires all to exist, as objects for His loving-kindness. Seeing then all created nature, as far as its own laws were concerned, to be fleeting and subject to dissolution, lest it should come to this and lest the Universe should be broken up again into nothingness, for this cause He made all things by His own eternal Word, and gave substantive existence to Creation, and moreover did not leave it to be tossed in a tempest in the course of its own nature, lest it should run the risk of once more dropping out of existence; but, because He is good He guides and settles the whole Creation by His own Word, Who is Himself also God, that by the governance and providence and ordering action of the Word, Creation may have light, and be enabled to abide alway securely.\endnote{St.\ Athanasius of Alexandria, \emph{Contra Gentes}, \S41.2--3 (Nicene and Post-Nicene Fathers, ser.\ 2, vol.\ 4, Robertson translation).}
\end{quote}

In the immediately following chapter, Athanasius unfolds that sustaining presence in cosmic music: the Word unites with the universe, ``quickening and sustaining all things everywhere, each severally and all collectively,'' holding mutually opposed elements---heat and cold, wet and dry---in concordant harmony, ``just as though some musician, having tuned a lyre\dots\ were to produce a single tune as the result, so also the Wisdom of God, handling the Universe as a lyre\dots\ produces well and fittingly, as the result, the unity of the universe and of its order, Himself remaining unmoved with the Father while He moves all things by His organising action, as seems good for each to His own Father.''\endnote{Athanasius, \emph{Contra Gentes}, \S42.1, 3.}

In the Latin West, St.\ Augustine of Hippo joined this theology of conservation to both anti-pantheistic precision and direct prayer. In his literal commentary on Genesis, \emph{De Genesi ad litteram} IV.12.22--23, Augustine reconciles Genesis's statement that God rested on the seventh day with Christ's declaration in John~5:17 that his Father is working even until now. God ceased creating new kinds of creatures, Augustine writes, yet he never ceased holding in being the creatures he had made. Here stands the very builder-and-house contrast that Aquinas would later adopt: \emph{Neque enim, sicut structor aedium cum fabricaverit, abscedit, atque illo cessante atque abscedente stat opus ejus; ita mundus vel ictu oculi stare poterit, si ei Deus regimen sui subtraxerit}---for not as a builder of houses walks away when he has finished constructing, and his work stands while he ceases and departs, could the world stand even for the blink of an eye if God withdrew his governance from it. Quoting Acts~17:28, Augustine adds an explicit anti-pantheistic safeguard: we are in God not as parts of his substance, since we are something other than he (\emph{cum aliud simus quam ipse}), yet solely through his unceasing work by which he holds all things together; hence if he withdrew this work of his from things, we would neither live, nor move, nor be (\emph{si hoc opus suum rebus subtraxerit, nec vivemus, nec movebimur, nec erimus}).\endnote{St.\ Augustine of Hippo, \emph{De Genesi ad litteram} IV.12.22--23 (Migne, PL~34); English renderings here translate the Latin text directly. Aquinas quotes this passage in \emph{ST} I, q.\ 104, a.\ 1, corpus.}

In the \emph{City of God}, Augustine ties continuous conservation to divine simplicity and the revealed name of Exodus~3:14. God alone is simple, he argues in Book~XI, chapter~10, ``because it has not anything which it can lose, and because it is not one thing and its contents another, as a cup and the liquor, or a body and its colour, or the air and the light or heat of it, or a mind and its wisdom. For none of these is what it has.'' In Book~XII, chapter~2, he adds that because God ``supremely is, and is therefore unchangeable, the things that He made He empowered to be, but not to be supremely like Himself,'' arranging natures in ranks of fuller or more limited existence so that no nature is contrary to God save non-entity itself. And in Book~XII, chapter~25, he restates the universal presence of the Creator's sustaining power:

\begin{quote}
It is His occult power which pervades all things, and is present in all without being contaminated, which gives being to all that is, and modifies and limits its existence; so that without Him it would not be thus or thus, nor would have any being at all\dots\ if He were, so to speak, to withdraw from created things His creative power, they would straightway relapse into the nothingness in which they were before they were created.\endnote{Augustine, \emph{De Civitate Dei} XI.10; XII.2; XII.25 (Marcus Dods translation, 1871).}
\end{quote}

When Augustine turns from cosmic exposition to direct prayer in \emph{Confessions} III.6.11, looking back on the years he wasted imagining God as a vast luminous body diffused through space, he compresses transcendence and innermost immanence into a single line: ``Thou wert more inward to me than my most inward part; and higher than my highest.''\endnote{Augustine, \emph{Confessions} III.6.11 (J.\ G.\ Pilkington translation, NPNF ser.\ 1, vol.\ 1).}

\subsection{Conciliar Definitions and Contemporary Catechetical Summary}

The Church has raised this core teaching on creation, divine transcendence, and continuous providence to solemn dogmatic definition.

Against Cathar and Albigensian dualism, which posited a second uncreated principle of material evil, the Fourth Lateran Council (1215) opened its profession of faith, \emph{Firmiter credimus}, by confessing one sole true God, eternal, immense, unchangeable, incomprehensible, almighty, and ineffable---\emph{unum universorum principium}, the single principle of all things, creator of all things visible and invisible, spiritual and corporeal, who by his almighty power at once from the beginning of time created out of nothing both orders of creatures, spiritual and corporeal, namely the angelic and the mundane, and then the human creature, constituted as it were in common out of spirit and body.\endnote{Fourth Lateran Council (1215), Const.\ 1, \emph{Firmiter credimus} (DS~800): \emph{unum universorum principium: creator omnium visibilium et invisibilium, spiritualium et corporalium: qui sua omnipotenti virtute simul ab initio temporis utramque de nihilo condidit creaturam, spiritualem et corporalem, angelicam videlicet et mundanam: ac deinde humanam, quasi communem ex spiritu et corpore constitutam}. The English rendering above translates this Latin text directly.}

Six and a half centuries later, confronting nineteenth-century pantheism, materialism, and rationalist emanationism, the First Vatican Council (1870) renewed and sharpened Lateran~IV in the Dogmatic Constitution \emph{Dei Filius}. In Chapter~1, \emph{Of God the Creator of All Things}, the council defines that the one true and living God, ``being one, sole, absolutely simple and immutable spiritual substance, is to be declared as really and essentially distinct from the world, of supreme beatitude in and from Himself, and ineffably exalted above all things beside Himself which exist or are conceivable.'' Immediately after reaffirming Lateran~IV on creation out of nothing by God's wholly free counsel, the council defines continuous providence: ``God protects and governs by His providence all things which He hath made, reaching from end to end mightily, and ordering all things sweetly'' (citing Wis.\ 8:1). Canon~5 of the same chapter seals the dogma with an anathema against anyone who refuses to confess ``that the world, and all things which are contained in it, both spiritual and material, have been, in their whole substance, produced by God out of nothing,'' or who claims that God created not by a will free from all necessity but as necessarily as he loves himself.\endnote{First Vatican Council (1870), Dogmatic Constitution \emph{Dei Filius}, ch.\ 1 and can.\ 5 \emph{De Deo rerum omnium creatore} (Latin text and Cardinal Manning's 1912 English translation).} While Lateran~IV and Vatican~I solemnly define creation out of nothing, the essential distinction between God and the world, and continuous providence, neither council canonizes Aquinas's internal real distinction of essence and \emph{esse} as binding dogma over Scotist or Su\'{a}rezian scholasticism.

Contemporary Catholic catechesis summarizes this conciliar and patristic inheritance directly. The \emph{Catechism of the Catholic Church} teaches that God stands infinitely above all his works yet remains present to the innermost depth of his creatures (CCC~300), and that having brought creation into being, God never abandons it to itself: he not only bestows being and existence at the beginning, but sustains and upholds every creature in being at each moment, enabling it to act and guiding it toward its end (CCC~301).\endnote{See \emph{Catechism of the Catholic Church}, nn.\ 300--301, summarized here doctrinally without verbatim reproduction.}
